\subsection{Complete results and physical ownership}
In the core runtime, dispatch binds session,
incarnation, job, input, target, and contract. A row returns either a
complete result with re-derived state/text or cancellation without partial
state. Complete rows can be installed and consumed before their batch
finishes. Only all terminal rows and device synchronization authorize
\texttt{BatchDrained} and executor reuse. The native process lifecycle
below has separate evidence.

\subsection{Covered work and safe-point exit: \texorpdfstring{$K$}{K}}
For the tested use model, $K$ requests exit when the installed Canonical
frontier covers a proposal's target and it has no remaining declared use.
This differs from $R$: valid work beyond the new root can remain useful.
Workers observe \texttt{CancelRequest} at decode boundaries; completion
can win, but each row has one terminal outcome. Static batches are not
compacted or refilled; masked costs and drain delays remain measured.

\subsection{Durable facts and outstanding input}
Before acknowledgment, synchronized frames bind parent/next states,
request identity, selected complete result, and receipt; recorded deltas
must reconstruct the state. Group-ready commit shares synchronization
among validated sessions without a collection timeout or cross-session
atomic transaction. Uncertain append/application outcomes block writes.
Write-ahead recovery~\cite{mohan1992aries} reconstructs facts and identities;
the native extension binds this to original correction calls.

\subsection{Validation boundaries}
Propositions~\ref{prop:rebase} and~\ref{prop:continuity} concern the stated
guards and compatible applications. The artifact retains an earlier,
separately frozen recovery prototype and its bounded model and crash
results; they do not validate the later core or native fault paths.
The selected-trace argument of Proposition~\ref{prop:refinement}, its
finite specification, and actual execution checks have distinct scopes.
